\scanpage{5}
6) Nécessité, en dim $>1$, d'une description plus délicate de la ``structure à l'infini'' d'un schéma lisse connexe \inserted{\unsure{séparé !}} sur $k$ alg. clos.

$X$
\diagram{\begin{tikzcd}[column sep=large]
1\arrow[r]&\pi_1^\infty(X)\arrow[r,"i"]&\pi_1(X)\arrow[r]&\pi_1(\widehat X)\arrow[r]&1\\
1\arrow[r]&\pi_1^\infty(U)\arrow[r]\arrow[u,"\alpha"]&\pi_1(U)\arrow[r]\arrow[u]&\pi_1(\mathfrak X)\arrow[r]\arrow[u]\arrow[d,equal]&1\\
&&&\pi_1(Y)&
\end{tikzcd}}
\inserted{\unsure{compactifié}}

\deleted{$\unsure{U\subset\widehat X\supset Y}$}
\diagram{\begin{tikzcd}
U\arrow[r,phantom,"\subset"]\arrow[d]&\mathfrak X\arrow[d]&\text{complété formel de}\;\widehat X\;\text{le long de }Y\\
X\arrow[r,phantom,"\subset"]&\widehat X\arrow[r,phantom,"\supset"]&Y(=\widehat X\setminus X)
\end{tikzcd}}

$\pi_1^\infty(X)$ est le sous-groupe \deleted{\unsure{clos}} \inserted{fermé} de $\pi_1(X)$ engendré par $\alpha(\pi_1^\infty(U))$, mais $\alpha$ n'est pas néc. surjectif.

Il faudrait expliciter, \inserted{en car. $0$ disons,} en termes de
\[
\Pi=\varprojlim\unsure{\pi_1}(U)
\]
et de la structure \inserted{``combinatoire''} de $Y$ (qu'on peut supposer diviseur à croisements normaux) le gpe $\pi_1^\infty(U)$ et l'action de Galois dessus, et exiger que \inserted{l'hom. ext.} $i\circ\alpha:\pi_1^\infty(U)\longrightarrow\pi_1(X)$ commute à l'action de Galois...

On doit pouvoir \inserted{par Lefschetz} faire une réduction au cas $\dim X=2$ ....

7) Sur l'espace de Teichmüller \inserted{complexe} $\widetilde M_{g,\nu\deleted{\unsure{\mathbb C}}}^{\unsure{an}}=\unsure{\mathcal M}_{g,\nu}^{\unsure{an}}$ (rev. universel de \unsure{la multiplicité} \inserted{\unsure{anal.}} $M_{g\nu}^{\unsure{an}}\unsure{(\mathbb C)}$)

\mspage{4}
